% ECONOMICS_PROBLEM: {"id": "O33-598", "title": "Adoption races and beliefs", "jel_code": "O33", "source_id": "OP-0413"}
% !TeX program = xelatex
\documentclass[11pt,a4paper]{article}
\usepackage[margin=23mm]{geometry}
\usepackage{fontspec}
\setmainfont{Latin Modern Roman}
\usepackage{amsmath,amssymb,mathtools}
\usepackage{xcolor,enumitem,booktabs,longtable,array,xurl,hyperref}
\definecolor{muted}{HTML}{53636C}
\hypersetup{colorlinks=true,urlcolor=blue,linkcolor=blue}
\setlength{\parindent}{0pt}
\setlength{\parskip}{5pt}
\newcommand{\R}{\mathbb R}
\newcommand{\E}{\mathbb E}
\newcommand{\N}{\mathbb N}
\newcommand{\ind}{\mathbf 1}
\newcommand{\Var}{\operatorname{Var}}
\newcommand{\Cov}{\operatorname{Cov}}
\newcommand{\supp}{\operatorname{supp}}
\DeclareMathOperator*{\argmin}{arg\,min}
\DeclareMathOperator*{\argmax}{arg\,max}
\newcommand{\entrymeta}[3]{{\sffamily\small\color{muted}\textbf{\texttt{#1}}\quad|\quad #2\par #3\par}}
\newcommand{\Problemref}[1]{\texttt{#1}}
\newcommand{\JELref}[2]{\texttt{#2} (JEL \texttt{#1})}
\begin{document}
\section*{Adoption races and beliefs}
\textbf{Problem ID:} \texttt{O33-598}\quad \textbf{JEL:} \texttt{O33}\par
% BEGIN ECONOMICS STATEMENT
\label{problem:OP-0413}
{\sffamily\small\color{muted}Original subject group: Innovation and AI economics\par}
\label{definitions:problem:30007695}
\textbf{Audit classification:} Background; proposed extension. The cited 2025 study was revised in August 2026 and already reports adoption-belief effects. Its measured intentions and robotics outcomes are not the catalogue proposed AI discovery welfare effect.
\subsection*{Definitions and precise research target}
Study firms in a specified industry that have not adopted a particular AI research tool at baseline. Firm \(i\) holds belief \(B_i\in[0,1]\) about the fraction of competitors adopting within a year, chooses research investment \(I_i\geq0\), and later produces validated innovation value \(Y_i\). Randomly provide accurate information \(Z_i\in\{0,1\}\) about competitors' previously measured adoption plans, without changing tool prices or access. Let \(B_i(z)\), \(I_i(z)\) and \(Y_i(z)\) denote potential beliefs and outcomes. Estimate \(E[B_i(1)-B_i(0)]\), \(E[I_i(1)-I_i(0)]\), and \(E[Y_i(1)-Y_i(0)]\) for eligible firms. Interpreting the investment ratio
\[
\beta=\frac{E[I_i(1)-I_i(0)]}{E[B_i(1)-B_i(0)]}
\]
as a belief response additionally requires a nonzero first stage and exclusion of other information channels; it is not automatically a structural derivative. Record actual competitor adoption to separate corrected expectations from subsequent strategic responses. To assess distortion, compare validated value net of real research costs under information provision and baseline, including spillovers within randomized market clusters. Does correcting adoption beliefs accelerate socially useful research or mainly induce duplication and defensive expenditure? This is a proposed extension of existing evidence on adoption expectations, with a defined information treatment and population.

\textbf{Definitions, assumptions and mathematical qualifications.} The signal treatment is a fixed truthful reporting rule about competitor adoption at a preregistered survey date. Potential outcomes are defined under that report and under no report, including each treatment version and access baseline. A market cluster rather than a single firm is the treatment unit if competitors react to others receiving information. The Wald ratio is defined only if \(E[B_i(1)-B_i(0)]\ne0\); its interpretation as a causal belief response needs a structural belief-to-investment map, instrument exclusion and an appropriate monotonicity or homogeneity restriction. A continuous belief change does not automatically produce a LATE interpretation. Social value includes rivals gains, duplication, real investment costs and any information-production costs.
\subsection*{Variants and assumption changes}
\begin{enumerate}
\item \textbf{Editorial, reduced form.} Estimate assignment effects on reported beliefs, verified actual adoption and validated outcomes; reduced-form effects do not require treating beliefs as an instrumented endogenous dose.
\item \textbf{Editorial, welfare channel.} Supply an identified belief-to-investment model and randomized market information saturation; assess net innovative value including duplication and strategic reactions.
\end{enumerate}
\subsection*{References and source audit}
\textbf{Checked source.} NBER Working Paper 34532, December 2025; revised August 2026, indexed abstract. The study reports experimental adoption-belief results, not the proposed innovation-welfare question. Primary metadata and abstract indexed; direct PDF returned HTTP 403. \url{https://www.nber.org/papers/w34532}.
\begin{enumerate}\raggedright
\item\hypertarget{definitions:source:30007695:1}{} Zoë B. Cullen; Ester Faia; Elisa Guglielminetti; Ricardo Perez-Truglia; Concetta Rondinelli. 2025. \emph{The Innovation Race: Experimental Evidence on Advanced Technologies}. NBER Working Paper 34532, December 2025; revised August 2026, indexed abstract.
\url{https://www.nber.org/papers/w34532}.
DOI: \url{https://doi.org/10.3386/w34532}.
\end{enumerate}

\subsection*{Common conventions: Source mathematical conventions}

These conventions apply to each entry unless explicitly replaced.

\textbf{Models and identification.} A primitive model \(m\) specifies all probability laws, feasible choices, information, equilibrium and measurement rules used by its target. The admissible class \(\mathcal M\) is fixed independently of outcomes. If \(P_O(m)\) is its observable probability law and \(\tau(m)\) the target, its identified set at observable law \(P\) is
\[
 \mathcal I_\tau(P)=\{\tau(m):m\in\mathcal M,\ P_O(m)=P\}.
\]
Point identification means that this set is a singleton. An empty set signals incompatibility of data and assumptions. Coordinate bounds need not describe the joint set. Bounds are sharp only if every retained target is attainable or arbitrarily approachable by compatible models. Finite bounds require explicit restrictions on unobserved quantities; unrestricted confounding, hidden wealth, extrapolation or arbitrary equilibrium selection can yield uninformative sets.

\textbf{Causal interventions.} A potential outcome \(Y_i(p)\) is the outcome under a complete policy regime \(p\), including timing, financing, information and market spillovers. The target population and units are fixed before treatment. With interference, use the complete assignment vector or a justified exposure mapping; individual no-interference notation alone cannot identify an economy-wide intervention. A difference of mean potential outcomes does not require the cross-world joint distribution. Conditioning on post-treatment migration, survival, enrollment or employment changes the estimand and needs separate justification.

\textbf{Statistical statements.} An estimator, confidence set, forecast or test specifies a sampling law, model class and loss/coverage criterion. Uniform \((1-\alpha)\)-coverage means \(\inf_{m\in\mathcal M}\Pr_m\{\tau(m)\in C_n\}\geq1-\alpha\), or an explicitly asymptotic version. The whole parameter space trivially has coverage, so informativeness requires a stated width, risk or power objective. A forecasting improvement is relative to a named benchmark, information set and evaluation population.

\textbf{Equilibrium and optimization.} An equilibrium comprises feasible plans, optimality, beliefs where needed, budget constraints and all market/resource clearing. The primitives or a stated benchmark define each correspondence \(\mathcal E(p;m)\). Nonemptiness is assumed or proved, never inferred just from compact policy sets. A selected-equilibrium target uses a declared selection rule; a robust target quantifies over all permitted equilibria. Use suprema or infima unless attainment is established. An \(\varepsilon\)-optimal policy is feasible and within \(\varepsilon\) of the finite optimum value. Report an empty feasible set separately.

\textbf{Welfare and accounting.} Normative utility, interpersonal normalization, social weights, numeraire, discounting and the use of public receipts are primitives. Behavioral utility need not coincide with the welfare criterion. Household welfare includes ownership income and rents once; adding profits again to compensating variation generally double-counts them. A monetary equivalent is defined by an explicit transfer and price convention, with monotonicity and range conditions. Average effects, mechanism contributions, transfers and welfare losses are different targets. Infinite sums require convergence and dynamic feasible plans require terminal or no-Ponzi conditions.

\textbf{Source versions.} A working-paper date, revision date and journal publication year are distinguished. A DOI for a journal version is not automatically the DOI of an earlier manuscript. Locators refer to the stated version and printed pages where specified. A metadata-only check does not verify a claim about a full-text conclusion. Every entry's source subsection narrows the provenance claim accordingly.
% END ECONOMICS STATEMENT
\end{document}
